\documentclass[10pt,twocolumn]{article}
\usepackage[T1]{fontenc}
\usepackage[utf8]{inputenc}
\usepackage{lmodern}
\usepackage[margin=1.8cm,top=2.4cm,bottom=2cm,columnsep=0.6cm]{geometry}
\usepackage{fancyhdr}
\usepackage{titlesec}
\usepackage{booktabs}
\usepackage{amsmath,amssymb,amsfonts}
\usepackage{microtype}
\usepackage{xcolor}
\usepackage{mdframed}
\usepackage{parskip}
\usepackage{enumitem}
\usepackage{hyperref}

\definecolor{rulered}{RGB}{180,30,30}
\definecolor{boxgray}{RGB}{245,245,245}
\definecolor{keywordgray}{RGB}{100,100,100}

\pagestyle{fancy}
\fancyhf{}
\fancyhead[L]{\textsc{\small Claude Mag}}
\fancyhead[C]{\small\textit{Machine Learning Research Digest}}
\fancyhead[R]{\small 2026-10-07}
\fancyfoot[C]{\small\thepage}
\renewcommand{\headrulewidth}{0.8pt}

\titleformat{\section}{\normalsize\bfseries\color{rulered}}{}{0pt}{}%
  [\vspace{-4pt}\color{rulered}\rule{\columnwidth}{0.4pt}]
\titlespacing{\section}{0pt}{8pt}{4pt}

\hypersetup{colorlinks=true,urlcolor=rulered,linkcolor=black}

\begin{document}
\setlength{\parindent}{0pt}
\setlength{\parskip}{4pt}

{\small\color{keywordgray}\textbf{Keywords:}
  video diffusion \textbullet{}
  physical consistency \textbullet{}
  seriality gap \textbullet{}
  autoregressive generation}\par\vspace{4pt}

{\LARGE\bfseries\raggedright
  S2PD: Serial-to-Parallel Diffusion for Physically and Logically
  Consistent Video Generation}\par\vspace{4pt}

{\large\itshape\raggedright
  A Two-Phase Diffusion Algorithm Resolves the Seriality Gap by Running
  Autoregressive Diffusion at High Noise and Parallel Diffusion at Low Noise}\par\vspace{6pt}

{\small\textbf{By} Jeffrey Hu, Daniel Olmeda Reino \& Ayush Tewari
  (Massachusetts Institute of Technology)}\par\vspace{2pt}

{\small\color{keywordgray}arXiv:2610.06847 \textbullet{} October 2026}
\par\vspace{2pt}\noindent{\color{rulered}\rule{\columnwidth}{0.6pt}}\vspace{6pt}

Bidirectional video diffusion models violate physical laws and logical rules
even when trained on unlimited in-distribution data---a structural obstacle
called the \emph{seriality gap}. S2PD resolves it with a single algorithmic
change: run the denoiser autoregressively at high noise, then switch to
parallel denoising at low noise, combining physical consistency with practical
sampling speed.

\begin{mdframed}[backgroundcolor=boxgray,linecolor=rulered,linewidth=1pt,
    innerleftmargin=6pt,innerrightmargin=6pt,
    innertopmargin=6pt,innerbottommargin=6pt]
\textbf{\small AT A GLANCE}\par\vspace{4pt}
\begin{description}[leftmargin=0pt,labelsep=4pt,itemsep=2pt,parsep=0pt,
    font=\small\bfseries]
  \item[Problem:] Bidirectional diffusion cannot reliably coordinate causally
    dependent video events, even with unlimited training data.
  \item[Why it matters:] Physics simulations, game environments, and embodied
    AI all require causally consistent video generation.
  \item[Method:] S2PD: autoregressive diffusion for the high-noise phase;
    parallel diffusion for the low-noise phase.
  \item[Result:] More reliable rule-following than bidirectional baselines;
    faster sampling than fully autoregressive methods.
  \item[Evidence:] Games, physical simulations, and real video benchmarks.
\end{description}
\end{mdframed}

\section{The Big Picture}

Video diffusion models have reached remarkable visual quality, but a persistent
blind spot---the \emph{seriality gap}---limits their utility for scientific
simulation, game generation, and embodied AI. Even when trained to convergence
on procedurally generated data with perfect ground-truth labels, bidirectional
diffusion models continue to violate physical laws: collision outcomes are wrong,
logical sequential rules are broken, and causal chains are severed.

Prior work (arXiv:2607.13031) identified and characterized this gap. S2PD is
the first method to close it without sacrificing the sampling speed advantage
of diffusion over fully autoregressive approaches.

\section{Why Existing Approaches Fall Short}

\textbf{More denoising steps do not help.} The seriality gap is not a capacity
issue. For deterministic video prediction, additional denoising steps do not
add effective serial computation---each step still processes all frames in
parallel.

\textbf{Architectural width does not help.} Wider networks also fail to close
the gap, confirming that the bottleneck is the absence of serial computation,
not model capacity.

\textbf{Fully autoregressive generation is slow.} Frame-by-frame autoregressive
generation provides the necessary serial compute but multiplies sampling time
by the number of frames, making it impractical for high-resolution or long
video generation.

\section{Core Mechanism}

The key insight is that different noise levels require different amounts of
serial computation:

\begin{itemize}[itemsep=2pt,parsep=0pt]
  \item \textbf{High noise} ($\sigma > \sigma^*$): the video is essentially pure
    noise; the denoiser must decide the \emph{global causal structure}---which
    events happen, in what order, with what consequences. This requires serial
    computation to propagate causal dependencies frame by frame.
  \item \textbf{Low noise} ($\sigma \leq \sigma^*$): the global structure is
    already established; the denoiser only needs to refine visual detail, which
    can be done in parallel without violating consistency.
\end{itemize}

S2PD switches between autoregressive and parallel denoising at threshold $\sigma^*$.

\section{Technical Formulation}

Let $\mathbf{x}_\sigma$ be the noisy video at noise level $\sigma$, with frames
$(x_{\sigma,1}, \ldots, x_{\sigma,T})$. The S2PD denoising trajectory is:

\textbf{Serial phase} ($\sigma > \sigma^*$): denoise frame by frame,
\[
\hat{x}_{0,t} = \epsilon_\theta(x_{\sigma,t},\, \sigma \mid x_{0,1:t-1}), \quad t = 1,\ldots,T
\]
where conditioning on previously denoised frames is implemented via masked
self-attention or cross-attention.

\textbf{Parallel phase} ($\sigma \leq \sigma^*$): denoise all frames jointly,
\[
(\hat{x}_{0,1}, \ldots, \hat{x}_{0,T}) = \epsilon_\theta(\mathbf{x}_\sigma, \sigma)
\]
starting from the latent at the end of the serial phase.

Total sampling time scales as approximately:
\[
\tau_\text{S2PD} \approx (T - 1)\,\tau_\text{serial-step} + \tau_\text{parallel}
\]
substantially less than $T\,\tau_\text{serial-step}$ for fully autoregressive
generation.

\section{Learning or Inference Procedure}

S2PD is a \textbf{training-free} inference algorithm applied to any pretrained
video diffusion model with autoregressive conditioning support.

\begin{enumerate}[itemsep=2pt,parsep=0pt]
  \item Sample Gaussian noise $\mathbf{x}_{\sigma_\text{max}}$.
  \item Run the DDIM solver serially (frame by frame) from $\sigma_\text{max}$
    to $\sigma^*$.
  \item Run the solver in parallel over all frames from $\sigma^*$ to $0$.
  \item Return the denoised video $\hat{\mathbf{x}}_0$.
\end{enumerate}

The switching threshold $\sigma^*$ is the only hyperparameter and is set
to the midpoint of the noise schedule in all experiments.

\section{What the Guarantee Says}

No formal consistency guarantee is provided. The empirical claim is: S2PD
produces fewer physical and logical violations than matched bidirectional
baselines across all tested environments, while requiring less sampling time
than fully autoregressive methods. The intuition is that the high-noise serial
phase establishes the causal skeleton, making the low-noise parallel phase
sufficient for detail refinement.

\section{Experimental Findings}

\begin{itemize}[itemsep=2pt,parsep=0pt]
  \item \textbf{Games:} S2PD follows Tetris collision rules and Sokoban push
    logic more reliably than bidirectional baselines.
  \item \textbf{Physical simulations:} Lower rule-violation rates on billiards
    and rigid-body interaction benchmarks.
  \item \textbf{Real video:} Improved temporal stability on real-world video
    prediction while maintaining visual quality.
  \item \textbf{Sampling efficiency:} Substantially faster than fully
    autoregressive generation while closing the consistency gap.
\end{itemize}

\section{Ablations and Interpretation}

\textbf{Switching threshold:} Earlier switching (less serial computation) degrades
consistency; later switching improves consistency but reduces efficiency gains.
The midpoint is Pareto-optimal.

\textbf{Phase ablations:} Removing the parallel phase (fully serial) matches
S2PD's consistency but increases sampling time by $T\times$; removing the
serial phase (fully parallel) restores the seriality gap.

\textbf{Error analysis:} Violations concentrate in frames requiring coordination
with distant prior frames, confirming that the gap is a long-range causal
dependency problem rather than a local temporal one.

\section*{Reference}

Jeffrey Hu, Daniel Olmeda Reino \& Ayush Tewari.
\textbf{S2PD: Serial-to-Parallel Diffusion for Physically and Logically
Consistent Video Generation}.
arXiv:2610.06847, October 2026.\\
\url{https://arxiv.org/abs/2610.06847}

\end{document}
